\part{Background}

\chapter{ECG Signal Processing Basics}
\label{chap:background}

Electrocardiogram (ECG) signals are 1D time-series recordings of cardiac electrical activity, characterized by key morphological features: P wave, QRS complex, and T wave. Arrhythmia classes are defined by the Association for the Advancement of Medical Instrumentation (AAMI): Normal (N), Supraventricular ectopic (S), Ventricular ectopic (V), Fusion (F), and Unknown (Q).

\section{Deep Learning for ECG Classification}
Prior work demonstrates:
\begin{itemize}
    \item 1D CNNs (ResNet) achieve high performance via local morphological feature extraction
    \item Bidirectional LSTMs capture long-term temporal dependencies in ECG rhythms
    \item Hybrid CNN-LSTM-Transformer models combine local, sequential, and global feature learning
    \item Multi-scale transformers (PhysioWave) simulate wavelet decomposition for noise-robust feature extraction
\end{itemize}

\section{Image-Based ECG Analysis}
Two advanced research directions are implemented:

\subsection{Image-to-Signal Pipeline}
Architecture: U-Net adapted for 2D-to-1D conversion (CNN encoder + 1D decoder). Process: Generate synthetic ECG images from 1D signals (simulating scanned medical records), train U-Net to reconstruct original signals. Innovation: Enables using 1D classifiers (ResNet, LSTM, etc.) on digitized ECG images.

\subsection{Image-to-Classification Pipeline}
Architecture: Vision Transformer (ViT) via timm library, with transfer learning from ImageNet. Process: Generate 2D ECG images from signals, directly classify using ViT without digitization. Innovation: End-to-end deep learning bypassing signal extraction, suitable for mobile photo diagnosis.

\section{Course Alignment}
This project aligns with the Deep Learning course (ESI, S4):
\begin{itemize}
    \item \textbf{CNNs:} 1D ResNet, multi-scale convolutions
    \item \textbf{RNNs:} Bidirectional LSTM
    \item \textbf{Transformers:} Multi-Module, PhysioWave
    \item \textbf{Transfer Learning:} Vision Transformer with ImageNet weights
    \item \textbf{Dimensionality Reduction:} PCA and t-SNE visualization
\end{itemize}
